\documentclass[10pt,a4paper]{article}
\usepackage[margin=18mm]{geometry}
\usepackage{amsmath,booktabs,array,tabularx}
\usepackage{cite,microtype}
\usepackage{xcolor}
\definecolor{SrcJin}{HTML}{56617B}
\definecolor{SrcWei}{HTML}{895B32}
\definecolor{SrcNS}{HTML}{317D37}
\definecolor{SrcHan16}{HTML}{7146A9}
\definecolor{SrcHan17}{HTML}{165CA7}
\definecolor{SrcCarey}{HTML}{947112}
\definecolor{SrcThesis}{HTML}{A5418A}
\definecolor{SrcWright}{HTML}{23828C}
\definecolor{Own}{HTML}{B53D36}
\definecolor{Standard}{HTML}{2A2A2A}
\newcommand{\Jin}[1]{\textcolor{SrcJin}{#1}}
\newcommand{\Wei}[1]{\textcolor{SrcWei}{#1}}
\newcommand{\NS}[1]{\textcolor{SrcNS}{#1}}
\newcommand{\Hsix}[1]{\textcolor{SrcHan16}{#1}}
\newcommand{\Hseven}[1]{\textcolor{SrcHan17}{#1}}
\newcommand{\Carey}[1]{\textcolor{SrcCarey}{#1}}
\newcommand{\Thesis}[1]{\textcolor{SrcThesis}{#1}}
\newcommand{\Wright}[1]{\textcolor{SrcWright}{#1}}
\newcommand{\Own}[1]{\textcolor{Own}{#1}}
\newcommand{\Std}[1]{\textcolor{Standard}{#1}}
\newcommand{\scite}[2]{\textcolor{#1}{\cite{#2}}}
\newcommand{\relB}{\textsf{B}}
\newcommand{\relA}{\textbf{\textsf{A}}}
\newcommand{\relD}{\textbf{\textsf{D}}}
\newcommand{\relS}{\textsf{S}}
\renewcommand{\citepunct}{], [}
\renewcommand{\citedash}{]--[}
\input{glyphtounicode}
\pdfgentounicode=1
\usepackage[hidelinks]{hyperref}
\hypersetup{pdftitle={An established traffic backbone for an elementary VSL task},
 pdfauthor={Differentiable traffic task review},
 pdfsubject={Literature comparison and formed modeling backbone; dependency hierarchy}}
\setlength{\parindent}{0pt}
\setlength{\parskip}{4pt}
\setlength{\abovedisplayskip}{4pt}
\setlength{\belowdisplayskip}{4pt}
\setlength{\abovedisplayshortskip}{3pt}
\setlength{\belowdisplayshortskip}{3pt}
\setlength{\tabcolsep}{4pt}
\renewcommand{\arraystretch}{1.10}
\newcommand{\head}[1]{\par\vspace{5pt}\noindent\textbf{#1}\par\nobreak\vspace{1pt}}
\newenvironment{points}{\begin{list}{$\bullet$}{%
 \setlength{\leftmargin}{1.3em}\setlength{\labelwidth}{0.6em}%
 \setlength{\labelsep}{0.4em}\setlength{\itemsep}{2pt}%
 \setlength{\parsep}{0pt}\setlength{\topsep}{2pt}}}{\end{list}}
\begin{document}
\begin{center}
{\large\bfseries An established traffic backbone for an elementary VSL task}\\[3pt]
Literature comparison and formed modeling backbone --- source hierarchy, 22 October 2026
\end{center}
\vspace{-3pt}

\head{1. What a suitable backbone must provide}
Can an established variable speed limit (VSL) model support the information,
control and gradient questions with few traffic-model changes? Consider an
ending lane E, its adjacent continuing lane A and farther lane B. Initial
states can share the same E and combined A+B traffic at every road position
while differing in A/B loading. We ask whether that difference changes merge
service, useful common-speed advice and its sensitivity. A suitable backbone
must represent the distinction, conserve traffic, respect receiving space
and count every initial unit's time until exit. Simplicity concerns the
traffic rules we must author and justify, not citation count.

\head{2. What the closest published models offer}
Aggregate VSL models support questions about traffic location or history;
lane models additionally resolve A/B loading. The inspected starting points
offer the following tradeoff.

\begin{center}
\textbf{Table 1. Published starting points and their fit to the question}\par
\vspace{3pt}
\begin{tabularx}{\linewidth}{@{}>{\raggedright\arraybackslash}p{0.19\linewidth}
 >{\raggedright\arraybackslash}p{0.37\linewidth}
 >{\raggedright\arraybackslash}X@{}}
\toprule
Published model & What is already supplied & Direct finding for this review\\
\midrule
\Jin{Jin and Jin}~\scite{SrcJin}{jinjin}
& Spatial and link-average traffic; VSL entry control, capacity drop and
upstream holding queues.
& Simpler if the question concerns \emph{where} traffic lies along the
road. It changes the A/B information question.\\[3pt]
\Wei{Wei, Han and Wang}~\scite{SrcWei}{wei}
& Link traffic retaining entry/speed history, interacting speed cohorts
and capacity drop.
& A strong aggregate alternative, but no explicit A/B loading. Native
advice acts on entering traffic.\\[3pt]
\NS{Nagalur Subraveti et al.}~\scite{SrcNS}{ns}
& Successive lane-cell densities, lane-change incentives and a lane-drop
case; \Hsix{Han congestion lineage}~\scite{SrcHan16}{han16}.
& Selected lane backbone. Adopt \Hseven{Han's VSL sending restriction}~\scite{SrcHan17}{han17}
and explicitly replace movement realization.\\[3pt]
\Carey{Carey et al.}~\scite{SrcCarey}{carey}
& Multilane cell transmission, mandatory movements and receiving-allocation
alternatives.
& Supports conservative service and forward fallback. More VSL/lane-end
task closure would be needed as the main backbone.\\
\bottomrule
\end{tabularx}
\end{center}

\textbf{Few-dependency construction.} Use \NS{Subraveti's lane framework} as
the backbone and follow its adoption of \Hsix{Han's congestion approach},
completing and adapting that formulation for the \Hseven{VSL task}.
Replace movement realization with one \Carey{Carey-led adaptation};
\Wright{Wright} supports its allocation rule. The thesis clarifies the backbone.
Thus \emph{one principal framework with inherited budgets and one adapted
movement component} carries the traffic model. Our \Own{scenario and composition}
remain explicit. This concentrates independent additions within the inspected
literature; global minimality is not proved. Jin--Jin is simpler if the
information question changes to spatial traffic differences.

\textbf{Source hierarchy and relationship key.} Columns show architectural
roles; colors credit rules and references, not ownership of symbols.
\par\nobreak\vspace{2pt}
\begin{tabularx}{\linewidth}{@{}>{\raggedright\arraybackslash}X>{\raggedright\arraybackslash}X>{\raggedright\arraybackslash}X@{}}
\toprule
\textbf{Principal framework}&\textbf{Inherited budget lineage}&\textbf{Movement adaptation}\\
\NS{Subraveti article} \scite{SrcNS}{ns}&
\Hsix{Han 2016} \scite{SrcHan16}{han16} (adopted in [3])&
\Carey{Carey et al.} \scite{SrcCarey}{carey} (principal)\\
\Thesis{Subraveti thesis} \scite{SrcThesis}{thesis} (clarification)&
\Hseven{Han 2017} \scite{SrcHan17}{han17} (VSL extension)&
\Wright{Wright et al.} \scite{SrcWright}{wright} (supporting)\\
\multicolumn{3}{@{}l@{}}{Comparison only: \Jin{Jin--Jin} \scite{SrcJin}{jinjin};
\Wei{Wei et al.} \scite{SrcWei}{wei}. \Own{This study}: scenario, composition and task (\relD).}\\
\midrule
\multicolumn{3}{@{}p{\linewidth}@{}}{\relB: source form retained/re-expressed; \relA: adapted/composed from the
colored source; \Own{\relD}: defined/derived here, including member choices
(not a novelty claim); \Std{\relS}: standard algebra, conservation or calculus.}\\
\bottomrule
\end{tabularx}\par

\head{3. The formed modeling backbone}
\textbf{3.1. Lane-cell traffic state.}
The \NS{lane-density architecture (\relB)} forms a
\emph{lane-resolved cell-transmission-type (CTM) model} with
\Hsix{inherited Han budgets}. Physical cells retain density and finite
storage throughout the road; there is no separate link transmission model
(LTM) link. E can change into A at upstream cell boundaries; traffic in the
final E cell must merge at the exit.
Unserved mass stays in its current cell, so congestion can extend upstream.
The causal chain is:
\begin{center}
\fbox{\NS{Lane stocks}}\ $\longrightarrow$\ \fbox{\NS{Desires} / \Hsix{budgets}}\
$\longrightarrow$\ \fbox{\Carey{Adapted movements}}\ $\longrightarrow$\ \fbox{\Std{New stocks}}
\end{center}
At step $k$, density $\rho^k_{i\ell}$ and continuous stock
$n^k_{i\ell}=\Delta x\rho^k_{i\ell}$ describe section $i$, lane
$\ell\in\{E,A,B\}$. These labels name current lanes: early-merged E traffic
joins A's stock. \Own{The declared member (\relD)} has $N=10$ sections,
$\Delta x=0.1$\,km and $\Delta t=1/600$\,h (6\,s).

\textbf{3.2. Core modeling assumptions.}
\Own{The task and model choices (\relD)} are:
\begin{points}
\item \textbf{Cohort and advice.} Initialize a finite cohort inside the road,
independently of advice; no later traffic enters. Choose one common
$v\in[30,60]$\,km/h at the start and hold it constant.
\item \textbf{Common lane physics.} Fix the geometry and use one common
nominal triangular flow--density diagram:
nominal free speed $u_0=60$, backward-wave speed magnitude $w=20$\,km/h,
jam density $k_J=150$\,veh/km/lane and maximum capacity-drop fraction
$\delta=0.1$. Initial densities lie in $[0,k_J]$; mass is continuous.
\item \textbf{Exit and accounting.} Trips end when traffic leaves through A
or B; downstream congestion does not block the exit. Count every initial
unit's time until it leaves, including traffic still queued at the lane end.
\item \textbf{Service and updates.} All movements use old traffic and space
and advance at most one section per step. Use equal positive priorities
and no extra merging-headway factor. The grid satisfies
$u_0\Delta t,w\Delta t\leq\Delta x$.
\item \textbf{Sensitivity.} Differentiate the complete objective where the
active traffic rules remain unchanged; switching cases need separate treatment.
\end{points}

\textbf{3.3. Sending and receiving capacity.}
Movement needs both traffic able to leave and room to accept it. The
\emph{sending rate} $D$ is how much a cell can release per hour; the
\emph{receiving rate} $R$ is how much the target lane can accept per hour
from all competing streams. Both are in vehicle-equivalents per hour.
\NS{Subraveti adopts Han for capacity drop}~\scite{SrcNS}{ns}; we follow that
\Hsix{inherited formulation (\relB)}~\scite{SrcHan16}{han16} with
\Hseven{nominal/VSL caps (\relB)}~\scite{SrcHan17}{han17}.
\Own{Our completion (\relD)} keeps the full sending ceiling, without turning
it off merely because the next cell becomes denser. This exact variant remains ours.
\Std{FD identities are standard (\relS).} The discharge ceiling $Q(x)$ equals nominal
capacity $C$ below critical density $\rho_c$, then falls linearly to
$(1-\delta)C$ at jam density $k_J$:
\begin{align}
 C&=\frac{u_0wk_J}{u_0+w},\quad \rho_c=C/u_0,\quad
 Q(x)=\min\!\left\{\Hseven{C},\Hsix{C\left[1-\delta\frac{x-\rho_c}{k_J-\rho_c}\right]}\right\},\label{eq:ceiling}\\
 \Hseven{D_{i\ell}}&=\Hseven{\min\{v\rho_{i\ell},Q(\rho_{i-1,\ell})\}},\quad
 \Hsix{R(x,y)=\min\{Q(x),w(k_J-y),w(k_J-x)+w_d(x-y)\}}.\label{eq:budgets}
\end{align}
Here $x,y$ are the current upstream/downstream densities in the receiving
lane. The three receiving bounds express discharge capacity, downstream
space and Han's additional congestion-wave constraint. Its fixed recovery
slope is \Hsix{$w_d=Q_J/(k_J-Q_J/u_0)$, $Q_J=(1-\delta)C$}.
Our \Hsix{lane-wise budget mapping (\relA)} gives $d_{i\ell}=\Delta t D_{i\ell}$ and
$r_{ij}=\Delta t R(\rho_{ij},\rho_{i+1,j})$, amounts available in one step.
All straight and diagonal
arrivals to target $j$ share that \emph{one} $r_{ij}$.
At the exit, $r_{Nj}=\Delta t R(\rho_{Nj},0)$, $j=A,B$, represents an
unblocked downstream road; upstream boundary density is zero. Advice caps sending;
receiving changes through the evolving densities, not directly through $v$.

\textbf{3.4. Lane-change incentives.}
This component describes why traffic wants to move over before the lane ends.
\NS{Subraveti's desire rule (\relB)}~\scite{SrcNS}{ns} compares perceived densities
$K$, with \Thesis{current/two-cell lookahead, weights $2{:}2{:}1$ (\relB)}~\scite{SrcThesis}{thesis}.
The dimensionless multiplier $I$ expresses the incentive for a particular
lane change; $P_{i,\ell j}$ is the desired lateral proportion before shared
traffic/service limits are applied:
\begin{equation}
 \color{SrcNS}
 P_{i,\ell j}=\left[\frac{I_{i,\ell j}K_{i\ell}-K_{ij}}
 {K_{i\ell}+K_{ij}}\right]_+,\qquad
 I=1+I^r+I^{kr}+I^{coop},\qquad [x]_+=\max(x,0).\label{eq:desire}
\end{equation}
With $I=1$, only a less crowded destination attracts traffic. A larger
multiplier can motivate a change even toward a more crowded lane. The added
terms represent route urgency, keep-right preference and cooperation.
Within the last $D_r=0.75$\,km,
\NS{E-to-A route bias (\relB)} grows as \NS{$I^r=(1-g_i/D_r)^3$}, where
$g_i=(N-i)\Delta x$; it is zero farther upstream. This adds urgency, not
a fixed merge fraction or a hard start line: density incentives can cause
earlier E-to-A movement. \Own{Member conventions (\relD)} select mode/direction
branches for keep-right/cooperation; free A can be encouraged toward B to
make room. Precise mode and finite-end conventions are in the companion
formulation; they add no separate traffic component.

\textbf{3.5. Movement allocation and forced merge.}
Desire does not guarantee a lane change: requests must share the available
traffic and receiving space. Before E ends, a failed attempt may still
continue forward in E; at its end, traffic must merge or wait.
\Carey{Movement adaptation (\relA)}~\scite{SrcCarey}{carey} replaces native
same-cell lateral/adjusted-forward coupling. \Own{Our composition (\relD)}
uses lane indices for each section: $\ell$ is the origin, $j$ the destination.
The desired lateral proportion $P_{\ell j}$ becomes an assigned fraction
$b_{\ell j}$ of the origin's sending amount $d_\ell$; $m_{\ell j}$ is the
resulting request. \Carey{Proportional normalization (\relA)} preserves desires whose sum
is at most one, otherwise scales them to fit one shared sending amount.
The remainder goes straight. Two 70\% lateral requests therefore become
50\% each, rather than spending 140\% of the same traffic:
\begin{equation}
 \Carey{b_{\ell j}=\frac{P_{\ell j}}{\max(1,\sum_{h\ne\ell}P_{\ell h})}\ (j\ne\ell)},
 \quad \Std{b_{\ell\ell}=1-\sum_{j\ne\ell}b_{\ell j}},\quad
 \NS{m_{\ell j}=d_\ell b_{\ell j}}.\label{eq:requests}
\end{equation}
We adapt \Carey{R2's proportional receiving rule (\relA)}~\scite{SrcCarey}{carey}
to the sender's budget; this placement is ours. The \Std{straight remainder (\relS)}
is accounting; \NS{request formation (\relB)} follows
\NS{\cite[Eq.~(14)]{ns}}.
\Carey{R3/A3 sharing and forward fallback (\relA)}~\scite{SrcCarey}{carey}
underpin the replacement. We extend its two-stream sharing to all target
requests; \Wright{Wright's priority formulation}~\scite{SrcWright}{wright}
provides supporting justification, without importing its general junction model.

\begin{samepage}
For each target, choose the common service level $\lambda_j$ so that the
served amounts $z_{\ell j}$ fill its receiving amount $r_j$, or satisfy all
requests if they are smaller:
\begin{equation}
 \color{SrcCarey}
 z_{\ell j}=\min(m_{\ell j},\lambda_j),\qquad
 \sum_\ell z_{\ell j}=\min\!\left(r_j,\sum_\ell m_{\ell j}\right).\label{eq:sharing}
\end{equation}
\end{samepage}
Small requests are fully served and release unused entitlement to the others.
\Carey{Forward continuation (\relA)} lets traffic not served in that allocation
still advance in its own lane if space remains. This additional flow
$y_\ell$ is limited by unused sending and unused same-lane receiving.
If receiving space also prevents forward movement, the traffic waits in
its current cell. Thus
\Own{$f_{\ell j}=z_{\ell j}+\mathbf1_{\{j=\ell\}}y_\ell$ (\relD)} is realized movement.
An unsuccessful E-to-A attempt can therefore continue in E if space remains.
\Own{Final-E rule (\relD): $b_{EA}=1$, $y_E=0$}; E must use A's outlet or wait.
It competes with A and any B-to-A request for the same outlet budget.
Equal priorities are not a constant realized merge proportion; demand and
receiving state determine the flow. At the choke, E uses part of A's exit
capacity, so A's own traffic may have to wait. The resulting accumulation
in A can reduce space for later arrivals and propagate congestion upstream.
There is no separate speed-mismatch or acceleration-recovery meter.
Our normalization and three-lane sequencing remain behavioral adaptations
beyond the lane-end geometry.

\Carey{Diagonal transport (\relA)} advances one section when changing lanes.
Apply old-state flows simultaneously using \Std{standard conservation (\relS)}:
\begin{equation}
 n^{k+1}_{i\ell}=n^k_{i\ell}-\sum_j f^k_{i,\ell j}
                             +\sum_h f^k_{i-1,h\ell}.\label{eq:update}
\end{equation}
Set $f_0=0$; $f_N$ counts traffic leaving the modeled road. Matched additions/subtractions
conserve mass, while donor/receiver limits protect stock and storage.

\textbf{3.6. Total system time and sensitivity.}
Minimize total system time $J$: all initial traffic's time until exit.
Once remaining mass is below
$\Delta x\rho_c$, all requests receive service and a unit in section $i$
has remaining mean time $(N-i+1)\Delta x/v$. Choose a fixed index $K_*$
inside that regime. The weight $h_{i\ell}=N-i+1$ counts the remaining
cell advances from section $i$ to the exit; $\mathbf1^\top n^k$ is the
mass still on the road. The complete objective is
\begin{equation}
 J(v)=\Delta t\sum_{k<K_*}\mathbf1^\top n^k
             +\Own{\frac{\Delta x}{v}h^\top n^{K_*}}.\label{eq:objective}
\end{equation}
\Std{Population-time accounting is standard (\relS)}; \Own{the exact residual is derived here (\relD)}.
The first term adds the time already spent by all unfinished traffic before
the cutoff. The second adds all travel time still to come: each remaining
stock is weighted by its distance to the exit and divided by advice speed.
Traffic remaining at the cutoff is therefore counted rather than discarded.
Write the complete update as $n^{k+1}=F(n^k,v)$. Between branch switches,
\Std{$s^0=0$, $s^{k+1}=F_n^{[k]}s^k+F_v^{[k]}$ (\relS)} propagates
$s^k=\partial n^k/\partial v$ through budgets, desires, allocation and stock.
Here $F_n,F_v$ are the update's first derivatives with respect to state and
advice. Then
\begin{equation}
 J'(v)=\Delta t\sum_{k<K_*}\mathbf1^\top s^k
       +\Own{\frac{\Delta x}{v}h^\top s^{K_*}
       -\frac{\Delta x}{v^2}h^\top n^{K_*}}.\label{eq:sensitivity}
\end{equation}
\Own{This model's full derivative (\relD)} uses \Std{standard chain rule (\relS)}.
Terms measure earlier traffic changes, residual-stock changes and
direct residual travel time. Hold a valid $K_*$ fixed locally; at native
mode or active-constraint switches a two-sided derivative can fail.

\textbf{3.7. Validation design and rival models.}
\Own{The validation design (\relD)} asks: does hiding A/B loading worsen
advice, can one local reading suffice, and do gradients help
choose advice? The planned comparisons keep the physical model and complete
objective fixed:

\emph{1. Isolate information value.} Compare an observer seeing E and A+B
at each section with one seeing each lane. For paired states with identical
reduced observations but different A/B loading, the reduced observer must
choose one shared advice; the full observer may choose separately.
Allow model-aware inference from admitted histories; decide before later
discharge reveals the split.

\emph{2. Test a compact information rival.} Add only final A-cell density
to E and A+B. Its value, 60 or 90\,veh/km/lane, identifies the reviewed
initial state. Compare with full information to test whether one reading
suffices; sufficiency beyond that pair is unproved.

\emph{3. Check gradients, then their usefulness.} First check Section 3.6's
complete-objective derivative against finite differences away from rule
switches. Then, within each information view and with matched initialization
and computational budget, compare gradient-guided advice with direct speed
search and maximum advice. Search is a strong scalar-action reference,
not a certified global optimum; maximum advice tests benefit over simply
using the highest speed. This separates missing information from poor
optimization. These checks are prospective, not reported results.
The \href{vsl_joined_candidate_contract_and_qualification.pdf}{companion formulation}
specifies bias branches and finite-domain conventions.

\par\pdfsavepos
\write16{REVIEW-MAIN-END-PAGE: \thepage; Y-SP: \the\pdflastypos}
\clearpage
\makeatletter
\renewcommand{\section}{\@ifstar{\head}{\head}}
\renewcommand{\@biblabel}[1]{\ifcase#1\or\Jin{[#1]}\or\Wei{[#1]}\or\NS{[#1]}\or\Hsix{[#1]}\or\Hseven{[#1]}\or\Carey{[#1]}\or\Thesis{[#1]}\or\Wright{[#1]}\else[#1]\fi}
\makeatother
\begin{thebibliography}{8}
\setlength{\itemsep}{4pt}\setlength{\parskip}{0pt}
\bibitem{jinjin} \Jin{H.-Y. Jin and W.-L. Jin, ``Control of a lane-drop bottleneck
through variable speed limits,'' arXiv:1310.2658, 2013 (archived manuscript
dated Oct. 29, 2018). [Online]. Available:
\url{https://arxiv.org/abs/1310.2658}.}
\bibitem{wei} \Wei{L. Wei, Y. Han, and M. Wang, ``A first-order link-based flow
model with variable speed limits and capacity drops for freeway networks,''
\emph{Comput.-Aided Civ. Infrastruct. Eng.}, vol.~40, no.~21,
pp.~3200--3217, 2025, doi:
\href{https://doi.org/10.1111/mice.13492}{10.1111/mice.13492}.}
\bibitem{ns} \NS{H.~H.~S. Nagalur Subraveti, V.~L. Knoop, and B. van Arem,
``First order multi-lane traffic flow model---an incentive based macroscopic
model to represent lane change dynamics,'' \emph{Transportmetrica B, Transp.
Dyn.}, vol.~7, no.~1, pp.~1758--1779, 2019 (archived article header: 2020),
doi: \href{https://doi.org/10.1080/21680566.2019.1700846}{10.1080/21680566.2019.1700846}.}
\bibitem{han16} \Hsix{Y. Han, Y. Yuan, A. Hegyi, and S.~P. Hoogendoorn,
``New extended discrete first-order model to reproduce propagation of jam
waves,'' \emph{Transp. Res. Rec.}, vol.~2560, pp.~108--118, 2016, doi:
\href{https://doi.org/10.3141/2560-12}{10.3141/2560-12}.}
\bibitem{han17} \Hseven{Y. Han, A. Hegyi, Y. Yuan, S. Hoogendoorn,
M. Papageorgiou, and C. Roncoli, ``Resolving freeway jam waves by discrete
first-order model-based predictive control of variable speed limits,''
\emph{Transp. Res. C, Emerg. Technol.}, vol.~77, pp.~405--420, 2017, doi:
\href{https://doi.org/10.1016/j.trc.2017.02.009}{10.1016/j.trc.2017.02.009}.}
\bibitem{carey} \Carey{M. Carey, C. Balijepalli, and D. Watling, ``Extending the
cell transmission model to multiple lanes and lane-changing,''
\emph{Netw. Spat. Econ.}, vol.~15, pp.~507--535, 2015, doi:
\href{https://doi.org/10.1007/s11067-013-9193-7}{10.1007/s11067-013-9193-7}.}
\bibitem{thesis} \Thesis{H.~H.~S. Nagalur Subraveti, \emph{Lane-specific Traffic
Flow Control}, Ph.D. dissertation, Delft Univ. Technol., Delft, The
Netherlands, 2021, ch.~4, doi:
\href{https://doi.org/10.4233/uuid:e405a549-b1fe-4fa1-8090-4429612359d9}{10.4233/uuid:e405a549-b1fe-4fa1-8090-4429612359d9}.}
\bibitem{wright} \Wright{M.~A. Wright, G. Gomes, R. Horowitz, and
A.~A. Kurzhanskiy, ``On node models for high-dimensional road networks,''
\emph{Transp. Res. B, Methodol.}, vol.~105, pp.~212--234, 2017, doi:
\href{https://doi.org/10.1016/j.trb.2017.09.001}{10.1016/j.trb.2017.09.001}.}
\end{thebibliography}
\end{document}
